% The method figure: coordinates are gauge -> ask in side branches -> each answer law on a quarter circle (Fisher-Rao =
% twice the arc) -> the geometry, compared with the hidden one; with full log-probabilities: the sphere -> exact Gram
\begin{figure*}[t]
\centering
\resizebox{\textwidth}{!}{%
\begin{tikzpicture}[font=\scriptsize, >={Stealth[length=4pt]},
  box/.style={draw=gbmute, rounded corners=2pt, line width=0.4pt, inner sep=3pt},
  ttl/.style={font=\scriptsize\bfseries, text=gbink, anchor=west},
  arr/.style={->, line width=0.6pt, draw=gbmute}]
% ---- panel 1: gauge
\begin{scope}[shift={(0,0)}]
  \node[ttl] at (-0.1,1.55) {(1) Coordinates are gauge};
  \draw[->, gbmute] (0,0) -- (1.1,0); \draw[->, gbmute] (0,0) -- (0,1.1);
  \begin{scope}[rotate=30]\draw[->, gbblue] (0,0) -- (1.1,0); \draw[->, gbblue] (0,0) -- (0,1.1);\end{scope}
  \fill[gborange] (0.75,0.45) circle (1.3pt); \fill[gborange] (0.3,0.85) circle (1.3pt);
  \draw[gborange, densely dotted] (0.75,0.45) -- (0.3,0.85);
  \node[anchor=west, text width=1.75cm, align=left] at (1.25,0.55) {$h$ and $Gh$ give the same outputs; distances and angles survive};
\end{scope}
% ---- panel 2: side branches
\begin{scope}[shift={(3.55,0)}]
  \node[ttl] at (-0.1,1.55) {(2) Ask in side branches};
  \node[box, fill=gbrow] (x) at (0.25,0.55) {$x$};
  \foreach \y/\q in {1.05/q_1, 0.55/q_2, 0.05/q_m}{
    \draw[arr] (x.east) -- (0.95,\y);
    \node[anchor=west] at (0.95,\y) {$\q$:};
  }
  \foreach \y/\p in {1.05/0.8, 0.55/0.25, 0.05/0.6}{
    \fill[gbblue!80] (1.42,\y-0.07) rectangle ++(\p*0.7,0.14);
    \draw[gbmute, line width=0.2pt] (1.42,\y-0.07) rectangle ++(0.7,0.14);
  }
  \node[anchor=west] at (0.95,-0.38) {$\Psi_j(x)=P(\textsc{yes}\mid x\oplus q_j)$};
\end{scope}
% ---- panel 3: quarter circle
\begin{scope}[shift={(7.25,0)}]
  \node[ttl] at (-0.1,1.55) {(3) Each answer law on a circle};
  \draw[gbmute, line width=0.4pt] (1.1,0) arc (0:90:1.1);
  \draw[->, gbmute] (0,0) -- (1.35,0) node[below, text=gbmute] {$\sqrt{1-p}$};
  \draw[->, gbmute] (0,0) -- (0,1.3) node[left, text=gbmute] {$\sqrt{p}$};
  \fill[gbblue] ({1.1*cos(70)},{1.1*sin(70)}) circle (1.4pt) node[above right=-1pt] {$x$};
  \fill[gborange] ({1.1*cos(25)},{1.1*sin(25)}) circle (1.4pt) node[right] {$x'$};
  \draw[gbink, line width=1.0pt] ({1.1*cos(25)},{1.1*sin(25)}) arc (25:70:1.1);
  \draw[gbmute, densely dotted] (0,0) -- ({1.1*cos(70)},{1.1*sin(70)}); \draw[gbmute, densely dotted] (0,0) -- ({1.1*cos(25)},{1.1*sin(25)});
  \node[anchor=west, text width=2.05cm, align=left] at (1.55,0.6) {Fisher--Rao $=2\times$ arc: $|\vartheta-\vartheta'|$, $\vartheta=2\arcsin\!\sqrt{p}$. Questions add as squares: exact, flat};
\end{scope}
% ---- panel 4: the geometry vs hidden
\begin{scope}[shift={(11.6,0)}]
  \node[ttl] at (-0.1,1.55) {(4) The observable geometry};
  \foreach \a/\b in {0.1/0.3, 0.4/0.9, 0.8/0.5, 1.15/1.05, 0.35/0.05}{\fill[gbblue] (\a,\b) circle (1.2pt);}
  \foreach \a/\b in {0.1/0.3}{\draw[gbmute, line width=0.3pt] (0.1,0.3) -- (0.4,0.9) -- (0.8,0.5) -- (0.1,0.3) (0.8,0.5) -- (1.15,1.05);}
  \node[anchor=west, text width=1.9cm, align=left] at (1.35,0.55) {distances $d_Q$; unique (\v{C}encov); vs hidden: Spearman, no fit};
\end{scope}
% ---- panel 5: exact recovery
\begin{scope}[shift={(15.1,0)}]
  \node[ttl] at (-0.1,1.55) {(5) With log-probabilities};
  \draw[gbmute, line width=0.4pt] (0.55,0.55) circle (0.5);
  \draw[gbmute, line width=0.3pt] (0.05,0.55) arc (180:360:0.5 and 0.15);
  \fill[gborange] (0.85,0.95) circle (1.2pt); \fill[gborange] (0.2,0.8) circle (1.2pt); \fill[gborange] (0.7,0.12) circle (1.2pt);
  \node[anchor=west, text width=1.45cm, align=left] at (1.1,0.55) {$\|u\|^2{=}d\Rightarrow b_i^{\top}Sb_i{=}d$: exact Gram from $\tfrac{d(d+1)}{2}$ contexts};
\end{scope}
\end{tikzpicture}}
\caption{\textbf{The method.} (1) Hidden coordinates are not observable (\cref{prop:gauge}). (2) Yes/no questions are asked in side branches; the real conversation never sees them. (3) A Bernoulli law $p$ is the point $(\sqrt{1-p},\sqrt p)$ on a quarter circle, and its Fisher--Rao distance to another is twice the arc between them; independent questions add as squares, so the angles are exact flat coordinates (\cref{thm:geometry}). (4) The distances between contexts are the observable geometry; we compare it with hidden geometry without fitting anything. (5) With full log-probabilities, the normalisation sphere gives a linear system for the exact Gram matrix of the final states (\cref{thm:exact}).}
\label{fig:method}
\end{figure*}
